\documentclass[oneside,12pt]{article}          % ----- article

\input{../../../Latex/inputs/loading_packages.tex}			% ----- Packages utiles
\input{../../../Latex/inputs/preambule_cours.tex}         	% ----- Préambule feuilles de TD
\input{../../../Latex/inputs/mymacros.tex}                 % ----- Macros mathématiques
\input{../../../Latex/inputs/mise_en_page.tex} 		   	% ----- Macros de mise en page

\fancyhead[R]{2\up{nd} générale}            % En-tête droit
\fancyhead[C]{Pourcentages et évolutions}
\fancyhf[CF]{\footnotesize{Lycée Polyvalent Irène et Frédéric Joliot-Curie}}           % En-tête central


%%%%%%%%%%%%%%%%%%%%%%%%%%%%%%%%%%%%%%%%%%%%%%%%%%%%%%%%%%%%%%%%%%%%%%%%%%%%%%%%%%%%%%%%%%%%%%%%%%%%%%%%%%%%%
%%%%%%%%%%%%%%%%%%%%%%%%%%%%%%%%%%%%%%%%%%%%%%%%%%%%%%%%%%%%%%%%%%%%%%%%%%%%%%%%%%%%%%%%%%%%%%%%%%%%%%%%%%%%%
%%%%%%%%%%%%%%%%%%%%%%%%%%%%%%%%%%%%%%%%%%%%%%%%%%%%%%%%%%%%%%%%%%%%%%%%%%%%%%%%%%%%%%%%%%%%%%%%%%%%%%%%%%%%%


\begin{document}

\section{Proportions et pourcentages}

\begin{rappel}{Proportion}
	\begin{minipage}{0.5\textwidth} \restoreskip
		Soit un ensemble A appelé \textbf{population} et un ensemble B inclus dans l'ensemble A appelé \textbf{sous-population}.
		
		La \textbf{proportion} de B par rapport à A est le nombre
		\[ p = \dfrac{\text{Nombre d'éléments de B}}{\text{Nombre d'éléments de A}} \]
	\end{minipage} \hspace{5mm}
	\begin{minipage}{0.45\textwidth}
		\begin{tikzpicture}
			\draw[line width = 2pt , color = newblue , fill = newblue!15]
				(2,0) ellipse (4 and 1.75) (4,0) node{A} ;
				
			\draw[line width = 2pt , color = newred , fill = newred!15]
				(1,0) ellipse (2 and 1.5) node{B} ;
		\end{tikzpicture}
	\end{minipage}
\end{rappel}

\begin{remarque}
	\begin{itemize}
		\item Une proportion est toujours un nombre compris entre $0$ et $1$.
		\item Plus une proportion est proche de $1$, plus la taille de la sous-population est proche de la taille de la population totale.
		\item On peut exprimer une proportion à l'aide de pourcentages en multipliant sa valeur par $100$, mais \textbf{les pourcentages ne sont pas des proportions.}
	\end{itemize}
\end{remarque}

\begin{exemple}{Calculs simples}
	Dans un club multi-sport comptant 500 adhérents, 200 d'entre eux pratiquent le tennis.
	
	On peut alors dire que la proportion des pratiquants de tennis par rapport au total des adhérents est : $p = \dfrac{200}{500} = \dfrac{2}{5} = 0,4$
	
	On peut aussi dire que $40\%$ des adhérents du club pratiquent le tennis.
\end{exemple}



\begin{propriete}{Proportion de proportion}
	\begin{minipage}{0.5\textwidth} \restoreskip
		Soit A une population, B une sous-population de A, et C une sous-population de B.
		
		On note $p_1$ la proportion de B  par rapport à A.\\
		On note $p_2$ la proportion de C par rapport à B.
		
		Alors la \textbf{proportion de C par rapport à A} vaut $p = p_1 \times p_2$
	\end{minipage}
	\begin{minipage}{0.5\textwidth}
		\begin{tikzpicture}
			\draw[line width = 2pt , color = newblue , fill = newblue!15]
				(2,0) ellipse (4 and 1.75) (4.5,0) node{A} ;
				
			\draw[line width = 2pt , color = newred , fill = newred!15]
				(1,0) ellipse (2.5 and 1.25) (-1,0) node{B} ;
				
			\draw[line width = 2pt , color = newgreen , fill = newgreen!15]
				(1.25,0) ellipse (1.5 and 0.75) (2,0) node{C} ;
		\end{tikzpicture}
	\end{minipage}
\end{propriete}

\begin{exemple}{Proportion de proportion}
	Dans la classe de Tijani, $60\%$ des élèves suivent la spécialité mathématiques, et parmi ces élèves $40\%$ suivent également la spécialité NSI.
	
	Pour déterminer la proportion des élèves de la classe suivant à la fois la spécialité mathématiques et la spécialité NSI, il faut multiplier les deux proportions. \\
	Ici, $p_1 = 0,6$ et $p_2 = 0,4$. \quad On obtient donc $p = 0,6 \times 0,4 = 0,24$.
	
	On conclut que $24\%$ des élèves de la classe de Tijani suivent les spécialités mathématiques et NSI.
\end{exemple}



\section{Variations de quantité}

\begin{definition}{Variation absolue}
	Lorsqu'une quantité évolue en passant d'une valeur initiale $V_I$ à une valeur finale $V_F$, alors on dit que sa \textbf{variation absolue} est $\Delta V = V_F - V_I$
\end{definition}

\begin{exemple}{Un cas pratique}
	Si un article coûtant au départ $160$\euro\, est soldé à $120$\euro, alors sa variation absolue est $ \Delta V = -40$\euro
\end{exemple}


\begin{definition}{Variation relative}
	Lorsqu'une quantité évolue en passant d'une valeur initiale $V_I$ à une valeur finale $V_F$, alors on dit que sa \textbf{variation relative} est $t = \dfrac{V_F - V_I}{V_I}$
	
	Ce nombre est parfois appelé \textbf{taux d'évolution}
\end{definition}

\begin{remarque}
	La variation relative permet de décrire une évolution en pourcentages. Par exemple si la variation relative est $t = 0,2$, alors on peut dire que la quantité a subi une augmentation de $20\%$
\end{remarque}

\begin{exemple}{Retour au soldes}
	Si l'on reprend l'exemple précédent, on peut calculer la variation relative :\\
	$t = \dfrac{160 - 120}{160} = \dfrac{-40}{160} = \dfrac{-1}{4} = -0,25$
	
	On peut alors conclure que l'article a été soldé à $-25\%$
\end{exemple}



\begin{definition}{Coefficient multiplicateur}
	Lorsqu'une quantité évolue en passant d'une valeur initiale $V_I$ à une valeur finale $V_F$, on appelle \textbf{coefficient multiplicateur} le nombre $C$ qui vérifie l'égalité $V_F = V_I \times C$
\end{definition}



\begin{propriete}{Coefficient multiplicateur et taux d'évolution}
	On considère une évolution dont la variation relative vaut $t$ et le coefficient multiplicateur vaut $C$
	\[ C = 1 + t \]
\end{propriete}

\begin{remarque}
	Augmenter de $x\%$ revient à multiplier par $C = 1 + \dfrac{x}{100}$ \\
	Diminuer de $x\%$ revient à multiplier par $C = 1 - \dfrac{x}{100}$
\end{remarque}



\section{Applications des coefficients multiplicateurs}

\subsection{Évolutions successives}

\begin{definition}{Évolutions successives}
	Lorsqu'une quantité évolue en passant d'une valeur $V_1$ à $V_2$ , puis de $V_2$ à $V_3$, puis de $V_3$ à $V_4$, etc. jusqu'à atteindre une certaine valeur $V_n$, on dit qu'elle a subi une \textbf{évolution globale} de $V_1$ à $V_n$
	
	Le coefficient multiplicateur associé à cette évolution globale est appelé le \textbf{coefficient multiplicateur global}.
	
	\begin{center}
		\begin{tikzpicture}
			% Les valeurs et coordonnées
			\node[draw = black , shape = ellipse , minimum width= 1.5cm , minimum height = 1cm , text = black , line width = 1.5pt]
				(V1) at (0,0) {$V_1$};
			\node[draw = black , shape = ellipse , minimum width= 1.5cm , minimum height = 1cm , text = black , line width = 1.5pt]
				(V2) at (3,0) {$V_2$};
			\node[draw = black , shape = ellipse , minimum width= 1.5cm , minimum height = 1cm , text = black , line width = 1.5pt]
				(V3) at (6,0) {$V_3$};
			\node[draw = black , shape = ellipse , minimum width= 1.5cm , minimum height = 1cm , text = black , line width = 1.5pt]
				(V4) at (9,0) {$V_4$};
			\node[draw = black , shape = ellipse , minimum width= 1.5cm , minimum height = 1cm , text = black , line width = 1.5pt]
				(Vn) at (12,0) {$V_n$};
			\draw (10.5,0) node{.....};
			
			% Flèches d'évolution
			\draw[-> , color = newblue , line width = 1.5pt , shorten <= 3pt, shorten >= 3pt]
				(V1.north) to[out = 60 , in = 120]
				node[pos=0.5, fill = lightgreen, inner sep = 1pt] {$\times C_1$}
				(V2.north) ;	
			\draw[-> , color = newblue , line width = 1.5pt , shorten <= 3pt, shorten >= 3pt]
				(V2.north) to[out = 60 , in = 120]
				node[pos=0.5, fill = lightgreen, inner sep = 1pt] {$\times C_2$}
				(V3.north) ;
			\draw[->, color=newblue, line width=1.5pt, shorten <=3pt, shorten >=3pt]
    			(V3.north) to[out=60, in=120]
    			node[pos=0.5, fill = lightgreen, inner sep = 1pt] {$\times C_3$}
    			(V4.north);
				
			% Flèche globale
			\draw[-> , color = newred , line width = 1.5pt , shorten <= 3pt, shorten >= 3pt]
				(V1.south) to[out = -40 , in = -140 , looseness = 0.5]
				node[pos = 0.5 , draw = none , above] {$\times C_{\text{global}}$} (Vn.south) ;	
		\end{tikzpicture}
	\end{center}
\end{definition}



\begin{propriete}{Coefficient multiplicateur global}
	Lors d'évolutions successives, si les coefficients multiplicateurs sont $C_1$ , $C_2$ , ... , $C_n$ , alors le coefficient multiplicateur global est donné par la formule :
	\[ C_{\text{global}} = C_1 \times C_2 \times \ldots \times C_n \]
\end{propriete}

\begin{exemple}{Application directe}
	Entre le mois de mars et le mois d'avril de cette année, un cinéma a perdu $10\%$ de ses spectateurs, mais entre avril et mai ils ont gagné $30\%$ de spectateurs. Quelle est l'évolution globale du nombre de spectateurs entre mars et mai ?
	
	On cherche à déterminer le coefficient global de ces évolutions successives, on va donc calculer les coefficients multiplicateurs de chaque évolution.
	
	Une diminution de $10\%$ correspond à un coefficient multiplicateur $C_1 = 1 - \dfrac{10}{100} = 0,9$
	
	Une augmentation de $30\%$ correspond à un coefficient multiplicateur $C_2 = 1 + \dfrac{30}{100} = 1,3$
	
	Le coefficient multiplicateur global est donc : $C_{\text{global}} = C_1 \times C_2 = 0,9 \times 1,3 = 1,17$
	
	On peut donc conclure que le cinéma a gagné $17\%$ de spectateurs entre les mois de mars et mai.
\end{exemple}

\begin{remarque}
	Comme on vient de le constater dans l'exemple ci-dessus, une diminution de $10\%$ suivie d'une augmentation de $30\%$ ne donnent pas une augmentation de $20\%$ comme on aurait pu le croire !
	
	\textbf{Il est interdit d'additionner des pourcentages lors d'évolutions successives}.
\end{remarque}



\subsection{Évolution réciproque}

\begin{definition}{Évolution réciproque}
	Lorsqu'une quantité évolue en passant d'une valeur initiale $V_I$ à une valeur finale $V_F$, on appelle \textbf{évolution réciproque} l'évolution qui permet de revenir de $V_F$ à $V_I$.
	
	Le coefficient multiplicateur associé à cette évolution réciproque est appelé \textbf{coefficient multiplicateur réciproque}.
\end{definition}



\begin{propriete}{Coefficient multiplicateur réciproque}
	\begin{minipage}{0.6\textwidth}
		Lors d'une évolution dont le coefficient multiplicateur est $C$, le coefficient multiplicateur réciproque est $C' = \dfrac{1}{C}$
	\end{minipage} \hspace{5mm}
	\begin{minipage}{0.3\textwidth}
		\begin{tikzpicture}
			% Les valeurs et coordonnées
			\node[draw = black , shape = ellipse , minimum width= 1.5cm , minimum height = 1cm , text = black , line width = 1.5pt]
				(VI) at (0,0) {$V_I$};
			\node[draw = black , shape = ellipse , minimum width= 1.5cm , minimum height = 1cm , text = black , line width = 1.5pt]
				(VF) at (4,0) {$V_F$};
			
			% Flèches d'évolution
			\draw[-> , color = newred , line width = 1.5pt , shorten <= 3pt, shorten >= 3pt]
				(VI.north) to[out = 60 , in = 120]
				node[pos=0.5, fill = lightblue, inner sep = 2pt] {$\times C$}
				(VF.north) ;	

			\draw[-> , color = newred , line width = 1.5pt , shorten <= 3pt, shorten >= 3pt]
				(VF.south) to[out = -120 , in = -60]
				node[pos=0.5, fill = lightblue, inner sep = 2pt] {$\times \dfrac{1}{C}$}
				(VI.south) ;
		\end{tikzpicture}
	\end{minipage}
\end{propriete}

\begin{exemple}{Cas d'application}
	Une usine produit des pièces pour moissonneuse. Le gérant a noté que le nombre de pièces défectueuses avait augmenté de $60\%$. S'il souhaite revenir au nombre départ, quelle évolution le nombre de pièces défectueuses doit-il subir ?
	
	On commence par déterminer le coefficient multiplicateur : $C = 1 + \dfrac{60}{100} = 1,6$
	
	On calcule le coefficient multiplicateur réciproque $C' = \dfrac{1}{C} = \dfrac{1}{1,6} = 0,625$
	
	Pour revenir au nombre de départ, il faudrait donc que le nombre de pièces défectueuses subisse une diminution de $37,5\%$
\end{exemple}

\begin{remarque}
	Comme on vient de le constater dans l'exemple ci-dessus, pour annuler une augmentation de $60\%$, il ne suffit pas d'appliquer une diminution de $60\%$ !
\end{remarque}

\end{document}